\subsection{Projective coefficient algebras over a perfect field}
\label{C:artinian-perfect-subsection}

Let $k$ be perfect of characteristic $p>2$, let $R$ be a
finite-dimensional commutative $k$-algebra, and put
$G_R=O_{m,k}\times_kR$ for the split group. We consider rational
$G_R$-representations that are finite projective over $R$,
with duality $M^*=\operatorname{Hom}_R(M,R)$.
The geometric canonical templates are those of
Theorem~\ref{C:artinian_algclosed_templates} and
Corollary~\ref{C:finite_algclosed_algebra_WG}.
The extension below classifies the indecomposable modules over $R$
and reduces their form coefficients to specified division algebras
with involution. The remaining arithmetic of those algebras and
forms is stated explicitly.

\begin{lemma}[Proved: projectivity and isomorphism reflection for finite coefficient algebras]
\label{C:finite-algebra-projectivity-descent}
Let $R$ be a finite-dimensional commutative algebra over any field
$k$, and let $k\subset L$ be a field extension.
For every finite $R$-module $M$, the module $M$ is projective over
$R$ if and only if $M_L$ is projective over $R_L$.
Hom spaces between finite-dimensional $R$-linear rational
representations commute with extension of the coefficient field.
If $L/k$ is algebraic, two such representations that become
isomorphic over $L$ are already isomorphic over $k$.
These assertions also hold within the full subcategory of finite
projective $R$-modules.
\end{lemma}
\begin{proof}
The first statement is the finite-dimensional case of faithfully
flat projectivity descent
\cite[Proposition 10.83.3, Tag 058S]{C-StacksArtinianPerfect}.
Here it has a direct linear proof. Choose a surjection
$\pi:R^N\to M$. Its domain and codomain are finite-dimensional
$k$-spaces. An $R$-linear splitting $s$ is specified by a finite
$k$-matrix satisfying the linear equations
$s(am)=a s(m)$ for the members of a fixed $k$-basis of $R$,
together with $\pi s=1$. If a solution exists after extending
scalars to $L$, this linear system over $k$ is consistent already
over $k$, so it has a solution there. This proves descent.
Ascent follows by tensoring the splitting.

Write the coactions of $M,N$ in a finite-dimensional subcoalgebra
of $k[O_m]$. Equivariance and commutation with a basis of $R$
are finitely many linear equations for a matrix $M\to N$.
Kernels of finite-dimensional linear systems commute with any
field extension, proving the Hom assertion.

For the isomorphism assertion, an isomorphism and its inverse
over an algebraic extension use coefficients in some finite
intermediate extension $L_0/k$. Forgetting its scalar structure
gives
$M^{\oplus[L_0:k]}\simeq N^{\oplus[L_0:k]}$.
The category here has finite-dimensional endomorphism algebras
and split idempotents, hence Krull--Schmidt decompositions with
local endomorphism rings. The same is true of the projective
subcategory: an idempotent summand remains projective over $R$,
and the Fitting decomposition is given by an idempotent.
Krull--Schmidt uniqueness cancels the common positive
multiplicity and gives $M\simeq N$.
This is the direct algebraic-extension argument of
Lemma~\ref{C:NoetherDeuring_rational}, with the $R$-action included.
\end{proof}

Write $\bar R=R\otimes_k\bar k=\prod_j\bar R_j$.
Each local factor is a finite-dimensional algebra over $\bar k$.
The geometric index set $\mathscr I_{\bar R}$ is the disjoint
union of the canonical indices over these factors, each module
extended by zero on the others. Set
$\Gamma=\operatorname{Gal}(\bar k/k)$.
Fix a global $k$-basis of $R$ and use its multiplication constants
to encode every full tuple before identifying its local factor.
For $\gamma\in\Gamma$, apply $\gamma$ to all coefficient matrices
of the full tuple encoding both the $R$-action and the coaction,
then canonicalize again. This defines an action on
$\mathscr I_{\bar R}$; it includes the permutation of the local
Artin factors. Fix the orders used in the geometric
canonicalization and choose the least member of each orbit.

\begin{theorem}[Proved: complete Galois indices over finite coefficient algebras]
\label{C:artinian_perfect_Galois_indices}
The resulting finite Galois orbits $\Omega$ are in bijection
with all indecomposable finite projective rational
$G_R$-representations $M_\Omega$.
There are specified division algebras
$D_\Omega=\operatorname{End}_{R,G}(M_\Omega)/\operatorname{rad}$
with centers $K_\Omega$ and positive integers $s_\Omega$ such that
\[
 (M_\Omega)_{\bar k}\simeq
       \bigoplus_{i\in\Omega}U_i^{\oplus s_\Omega},\qquad
 [K_\Omega:k]=|\Omega|,\qquad
 s_\Omega^2=\dim_{K_\Omega}D_\Omega .                       \tag{C.3.P1}
\]
There is no remaining module-isomorphism quotient in these
indices. The representative and its division algebra are obtained
from a finite-dimensional endomorphism algebra by the construction
in the proof. Coefficient degree is preserved by this
correspondence. Scalar extension embeds the split Green ring
of these projective representations into its geometric split
Green ring, with image the subgroup having basis
\[
                  s_\Omega\sum_{i\in\Omega}[U_i].          \tag{C.3.P2}
\]
Products, duals and exterior operations in this subgroup are
computed geometrically over $\bar R$, with coefficient division
by the corresponding $s_\Omega$.
\end{theorem}
\begin{proof}
An indecomposable projective $R$-representation $M$ has a
finite-dimensional local endomorphism algebra $E$.
Its radical $J$ is nilpotent and $D=E/J$ is a finite-dimensional
division algebra over $k$. The center $K$ is finite separable
over $k$ because $k$ is perfect. The splitting theorem for
central simple algebras
\cite[Lemma 11.8.6, Tag 0753]{C-StacksPerfectDescent} gives
\[
 D\otimes_k\bar k\simeq
       \prod_{\tau:K\hookrightarrow\bar k}M_s(\bar k),
                   \qquad s^2=\dim_KD .                   \tag{C.3.P3}
\]
By the Hom comparison in
Lemma~\ref{C:finite-algebra-projectivity-descent},
$E\otimes_k\bar k=\operatorname{End}_{\bar R,G}(M_{\bar k})$.
Its radical is $J\otimes_k\bar k$: this ideal is nilpotent
and its quotient (C.3.P3) is semisimple.
All summands in a geometric Krull--Schmidt decomposition are
projective over $\bar R$, hence belong to the classified index
set. An indecomposable one is supported on one local factor
and has residue endomorphism field $\bar k$.
Comparing the semisimple quotient with (C.3.P3) gives
$[K:k]$ distinct indices, each with multiplicity $s$.
The Galois action on the embeddings of $K$ is transitive, so
these indices form one orbit. This proves (C.3.P1) for any
indecomposable already defined over $R$.

Conversely fix $U_i$ in a geometric orbit.
Its coefficient matrices, and the finite idempotent presentation
of its underlying projective $\bar R$-module, involve only
finitely many elements of $\bar k$. Choose a finite Galois
extension $L/k$ containing them. The idempotent, coaction
identities and inclusion into a finite free $\bar R$-module
then define a projective representation $U$ over $R_L$ whose
extension is $U_i$. Alternatively, projectivity follows by
the reflection lemma. This verifies that finite descent
has stayed in the projective category.

Let $W=\operatorname{Res}_{R_L/R}U$.
Since $R_L$ is finite free over $R$ and $U$ is a summand of
a finite free $R_L$-module, $W$ is projective over $R$.
Its coaction is $R$-linear, obtained by replacing each
$L$-coefficient by its multiplication matrix in an $L/k$-basis.
Separability gives the actual projective-module decomposition
\[
 W_{\bar k}\simeq
   \bigoplus_{\sigma:L\hookrightarrow\bar k}{}^\sigma U.    \tag{C.3.P4}
\]
The support is exactly $\Omega$, including the conjugate
Artin factors, with the same multiplicity at every member.

If $M,N$ are indecomposable summands of $W$, the first part
shows that each has all of $\Omega$ as geometric support.
Let their common-on-orbit multiplicities be $s,t$.
Then $(M^{\oplus t})_{\bar k}\simeq
(N^{\oplus s})_{\bar k}$.
Isomorphism reflection and Krull--Schmidt uniqueness give
$M\simeq N$ and $s=t$.
The same comparison with any other indecomposable supported
on $\Omega$ proves uniqueness. Thus $W\simeq M_\Omega^{\oplus r}$.

Compute $E_W=\operatorname{End}_{R,G}(W)$ by the finite linear
equations for both actions. It has semisimple quotient
$M_r(D_\Omega)$. Lift a primitive idempotent across its
nilpotent radical and take its image on $W$.
This image is the promised projective representative
$M_\Omega$. Every such choice has the same module index.
In particular the construction makes no assumption that an
individual geometric module descends to its field of moduli.

The coefficient filtration is $\bar R\otimes_{\bar k}C_d$,
where $C_d$ is defined over the prime field.
Galois conjugacy, restriction of scalars and constant change
of basis preserve it. Passing to a direct summand cannot
increase the degree, and faithful field extension reflects
membership in a filtration piece. Hence the least coefficient
degrees in (C.3.P1) agree.

The split Green group has its Krull--Schmidt basis
$[M_\Omega]$. Their geometric images have disjoint supports,
proving injectivity and (C.3.P2). Here all tensor products and
exterior powers are taken over $R$ or $\bar R$, not over
$k$ or $\bar k$. They commute with scalar extension and
remain projective. The algorithms of
Corollary~\ref{C:finite_algclosed_algebra_WG} therefore expand
their geometric images. At any $j\in\Psi$ the resulting
coefficient is the coefficient of $M_\Psi$ multiplied by
$s_\Psi$; division is consequently integral and gives exactly
that coefficient. The $R$-dual of a projective module commutes
with field extension. Thus duality sends an orbit to its
geometric dual orbit and preserves its Schur index.
Virtual exterior operations follow from the direct-sum
identity coefficientwise.
\end{proof}

\begin{theorem}[Proved reduction: the exact remaining form arithmetic]
\label{C:artinian_perfect_WG_reduction}
Use the complete indices of
Theorem~\ref{C:artinian_perfect_Galois_indices}.
For each self-dual $\Omega$, choose an actual nonsingular
$\varepsilon_\Omega$-symmetric $R$-linear form
$q_\Omega$ on $M_\Omega$, where
$\varepsilon_\Omega\in\{1,-1\}$, and let $\sigma_\Omega$ be
the resulting adjoint involution on $D_\Omega$.
Choose one member of each other dual pair.
There is an exact additive decomposition
\[
 WG_{\mathrm{sp}}(BG_R)=
 \bigoplus_{\Omega=\Omega^*}
      GW^{\varepsilon_\Omega}(D_\Omega,\sigma_\Omega)
 \ \oplus\
 \bigoplus_{\{\Omega,\Omega^*\},\,\Omega\ne\Omega^*}
      \mathbb Z\,h_\Omega .                              \tag{C.3.P5}
\]
The groups on the right are the split groups of nonsingular
hermitian multiplicity forms, and
$h_\Omega=[H(M_\Omega)]$.
The algebras, involutions and representative matrices are
specified by finite-dimensional constructions. Products
and exterior operations are determined by actual tensors
and exterior powers over $R$, decomposition, and transfer
of the resulting multiplicity forms to these residue algebras.

This reduces the nonfield problem to the indicated arithmetic
of explicit division algebras with involution. It does not
classify hermitian forms over every arbitrary perfect field
by numerical invariants, or assert a uniform algorithm for
that arithmetic. When algorithms for the indicated
finite-dimensional algebras and hermitian forms are provided,
the construction computes every specified ring or exterior
coefficient in finitely many such operations.
\end{theorem}
\begin{proof}
The projective $R$-representation category is additive and
idempotent complete. Its endomorphism algebras are finite
dimensional over $k$; Fitting gives local endomorphism rings
on indecomposables, and every radical is nilpotent and hence
complete. The $R$-linear dual is an exact strong duality on
finite projectives. These verify all three categorical
hypotheses in the hermitian Krull--Schmidt theorem
\cite[Theorems 3.2--3.3, pp.272--273]{C-QSS1979}.
In particular this application makes no Gorenstein or
Frobenius assumption on the coefficient algebra $R$.

For a self-dual indecomposable, an equivariant
isomorphism $g:M\to M^*$ exists by definition.
Write $2g=(g+g^*)+(g-g^*)$.
After composing with $g^{-1}$, a sum of two nonunits in the
local endomorphism ring cannot be the unit $2$.
Hence one summand is an isomorphism. It gives the required
symmetric or alternating $q_\Omega$, precisely the elementary
argument in \cite[Proposition 2.5(1), p.271]{C-QSS1979}.
The adjoint preserves the radical, giving $\sigma_\Omega$.
The equations for the $R$-linear equivariant form and its sign
are linear over $k$; nonsingularity is the nonvanishing of
its determinant as a $k$-linear map to the $R$-dual.
For infinite $k$ a finite determinant grid finds a nonsingular
member whenever one exists, and the finite-field case permits
finite enumeration.

Transfer with $q_\Omega$ identifies forms on sums of
$M_\Omega$ with $\varepsilon_\Omega$-hermitian forms over
its endomorphism algebra
\cite[Proposition 2.4, p.271]{C-QSS1979}.
Reduction modulo the nilpotent radical classifies these forms
by the corresponding forms over $D_\Omega$
\cite[Theorem 2.2 and Corollary 2.3, pp.269--270]{C-QSS1979}.
The sign is $\varepsilon_\Omega$ because the original form
being classified is symmetric. On a non-self-dual pair,
the theorem makes every symmetric object hyperbolic and
its multiplicity is a nonnegative integer.
Orthogonal decomposition by these dual orbits is unique.
Taking group completions proves (C.3.P5), with no exact-GW
metabolic relations imposed.

To compute an operation, choose actual representatives of
the input multiplicity forms, lift them through the nilpotent
endomorphism radicals using the same reduction theorem, and
apply transfer to obtain actual symmetric representations.
Use their tensor Gram matrices, or their determinant Gram
matrices on exterior powers, over $R$.
The finite-dimensional endomorphism algebra of the output
provides a Krull--Schmidt decomposition. Its indices are
identified by geometric canonicalization and
Theorem~\ref{C:artinian_perfect_Galois_indices}.
Transfer and radical reduction recover the exact residual
hermitian forms, which are retained as coefficients in (C.3.P5).
This is where their arithmetic is genuinely needed: geometric
multiplicities alone do not recover these form coefficients.
All constructions on any specified input have finite
dimension, and the virtual exterior formula uses finitely
many coefficients for each requested degree. This proves
the operation rule and the stated conditional algorithmic scope.
\end{proof}

\begin{corollary}[Proved: all Artin rings of prime characteristic with perfect residues]
\label{C:perfect-residue-Artin-reduction}
The preceding indecomposable classification, split Green
calculation and hermitian-arithmetic reduction apply to every
commutative Artinian ring $R$ of characteristic $p>2$ whose
residue fields are perfect, whether or not a coefficient
field is initially specified. They apply separately to its
finite local factors. In particular the only retained inputs
are the finite algebra structure, the explicit geometric
indices, and the residual division-algebra and form arithmetic.
\end{corollary}
\begin{proof}
An Artinian ring is a finite product of local Artinian rings,
so consider one with maximal ideal $\mathfrak m$ and perfect
residue field $K$. A coefficient field exists by the Cohen
structure theorem
\cite[Theorem 10.160.8, Tag 032A]{C-StacksArtinianPerfect}.
Here the perfect prime-characteristic case has a direct
formula. Choose $N$ with $\mathfrak m^{p^N}=0$.
For $a\in K$, choose any lift $b$ of the unique
$p^N$th root $a^{1/p^N}$, and define
\[
                         s(a)=b^{p^N}.                   \tag{C.3.P6}
\]
The result is independent of the lift, since changing $b$
by $c\in\mathfrak m$ changes its $p^N$th power by $c^{p^N}=0$.
Frobenius additivity and multiplicativity show that $s$ is
a unital ring homomorphism, and its composite with the
residue map is the identity. Thus $s(K)$ is a coefficient
field. Any other coefficient-field section has this same
value on $a$, by applying it to $a^{1/p^N}$, so the section
is in fact unique.

The finite filtration by powers of $\mathfrak m$ has
finite-dimensional $K$-vector-space quotients.
Consequently $R$ is finite dimensional over its coefficient
field. Apply the preceding results with $k=K$, then take
the finite product over the local factors.
The construction of $s$ uses inverse Frobenius in the
perfect residue field; in an effective version that
operation and lifts of residue elements are part of the
specified field and ring access.
This corollary does not include rings of characteristic
$p^e$ with $e>1$, or imperfect residue fields.
\end{proof}

\begin{theorem}[Proved: a complete effective nonfield answer over perfect $C_1$ fields]
\label{C:artinian_effective_C1}
Let $k$ be a perfect $C_1$ field of characteristic $p>2$,
and let $R$ be a finite-dimensional commutative $k$-algebra.
The projective Galois indices above have Schur index one.
For each self-dual index choose $q_i,\varepsilon_i$ as in
Theorem~\ref{C:artinian_perfect_WG_reduction}, and write its
residue endomorphism field and involution as $(K_i,\sigma_i)$.
The complete additive answer is
\[
 \begin{split}
 WG_{\mathrm{sp}}(BG_R)
   &\simeq
    \bigoplus_{i\in\mathscr J_{\mathrm s}}
       \bigl(\mathbb Z\oplus K_i^\times/(K_i^\times)^2\bigr)\\
   &\quad\oplus
    \bigoplus_{i\in\mathscr J_{\mathrm a}\cup
                    \mathscr J_{\mathrm u}\cup
                    \mathscr J_{\mathrm h}}\mathbb Z .
 \end{split}                                               \tag{C.3.P7}
\]
Here the four types have the same meaning as in (C.3.Q2),
with $R$-duality throughout.

Assume in addition that $k$ has the effective data specified
before Lemma~\ref{C:C1_algebra_certificates}, and that $R$ is
given by a finite $k$-basis with multiplication constants.
Then all finite projective rational $G_R$-representations
have decidable, repetition-free natural-number matrix indices.
The indecomposable indices, duality, all fields and involutions
in (C.3.P7), canonical square-class coordinates, equality,
products and every specified exterior-power coefficient are
computed by terminating algorithms. The indices use the
underlying $k$-dimension and least coaction degree; they do
not require a free $R$-module or a constant rank on all factors.
There is no remaining representation or form isomorphism
quotient in this effective answer.

For finite $k$, the same conclusion follows instead from
Theorem~\ref{C:Om-finite-base-algorithm}, since $R$ is a finite
ring. In particular, the infinite nonclosed fields
$k_{p,\ell}$ of Corollary~\ref{C:C1_tower_example}, with any
finite-dimensional coefficient algebra over them, give
concrete nonfield families to which the ordinary terminating
algorithm applies.
\end{theorem}
\begin{proof}
Every finite extension of $k$ is $C_1$ and has trivial Brauer
group by Lemma~\ref{C:C1_field_forms}. Thus the division
algebras in (C.3.P1) are fields and all Schur indices are one.
Apply that lemma's complete symmetric, alternating and
unitary form classification to (C.3.P5). It gives exactly
(C.3.P7): rank and ordinary determinant in symmetric type,
half-rank in alternating type, rank in unitary type, and
one hyperbolic multiplicity for a non-self-dual pair.
This proves the structural formula for arbitrary perfect
$C_1$ fields before any effective hypothesis is imposed.

We now supply the additional nonfield steps in the algorithm.
Let $a_1=1,\ldots,a_c$ be the supplied basis of $R$, let
$Q_i$ be its regular multiplication matrices, and let
$C_d^*$ have its known finite multiplication table.
On $k^N$ enumerate all tuples $(B_1,\ldots,B_c,T_1,\ldots,T_r)$
satisfying the multiplication and unit relations of $R$
and $C_d^*$, and $B_iT_a=T_aB_i$.
The equations are the full tuple equations (C.3.A1), now
without a freeness condition. The entries have effective
natural-number codes. The minimal-degree condition is again
the vanishing test for coefficients outside $C_{d-1}$.

\emph{A decidable projectivity test.}
Let $M=k^N$ be the resulting $R$-module, with basis $v_j$.
The explicit surjection
\[
 \pi:R^N\longrightarrow M,\qquad
              (r_1,\ldots,r_N)\longmapsto\sum_j r_jv_j
                                                               \tag{C.3.P8}
\]
has, on the basis vector $a_i$ in its $j$th summand,
the known column $B_iv_j$.
An $R$-linear section $s:M\to R^N$ is exactly a solution
of the finite affine linear system
\[
       sB_i=(I_N\otimes Q_i)s\quad(1\leq i\leq c),
                  \qquad \pi s=I_N.                         \tag{C.3.P9}
\]
Gaussian elimination decides consistency. A solution exists
if and only if $M$ is projective. If one exists, the computed
$e=s\pi$ is an explicit projective idempotent.
Every block of $e$, in the regular $R$-coordinates, commutes
with the $Q_i$; its value on $1$ recovers its entry in $R$.
Thus the algorithm has both the $k$-tuple and an actual
finite projective presentation over $R$.

Discard the nonprojective tuples. Compute the joint
centralizer of the remaining tuple. It is the complete
finite-dimensional algebra $\operatorname{End}_{R,G}(M)$.
Lemma~\ref{C:C1_algebra_certificates} gives a terminating
positive-certificate computation of its radical and
semisimple quotient, and hence decides localness and
produces indecomposable summands. All summands remain
projective. The intertwiner space includes both the
$R$-action and the coaction. Its determinant polynomial
therefore decides isomorphism, and the finite grid
constructs an isomorphism when one exists.
As in Theorem~\ref{C:Om_effective_C1}, retaining the least
isomorphic code among the finitely many codes up to the
input's code gives exactly one canonical index.
All rational projective representations occur in this
enumeration, because their coactions have finite degree.

\emph{Computing the correct dual and form types.}
Compute
$M^*=\operatorname{Hom}_R(M,R)$ as the vector space of
matrices $F:k^N\to k^c$ satisfying
$FB_i=Q_iF$ for every $i$.
Multiplication on the values supplies the $R$-action
on this space. If $s_{ab}$ are the antipode constants,
the dual coefficient operator sends
\[
                       F\longmapsto
                             \sum_b s_{ab}F T_b .
                                                               \tag{C.3.P10}
\]
This is precomposition on $R$-valued functionals, rather
than replacement by the ordinary $k$-linear dual.
Projectivity gives the usual evaluation isomorphism
$M\simeq M^{**}$, and the algorithm computes it in these
bases. Canonicalizing this full tuple gives the dual index.

An $R$-bilinear invariant form on $M$ is equivalently an
equivariant $R$-linear map $M\to M^*$.
Its symmetry or alternating condition is the equality
of its two $R$-valued evaluations, with the appropriate
sign. These are finite linear equations over $k$.
The dimensions of $M$ and $M^*$ over $k$ are equal, as
is seen on each local projective rank stratum.
Nondegeneracy is exactly a nonzero determinant of this
$k$-linear map. A finite determinant grid therefore
chooses a nonsingular form of one of the two signs
for a self-dual index, whose existence was proved above.
Adjoints are computed by composing this isomorphism,
the explicit $R$-dual maps and its inverse.
The certificate lemma computes the residue field
$K_i$ and its adjoint involution.

Lemma~\ref{C:C1_field_forms} now supplies complete
form coordinates. Square-class equality and its least
representative are computed by factorization over
$K_i$ and finite predecessor search, exactly as in
Theorem~\ref{C:Om_effective_C1}. Actual coordinate
representatives are obtained by lifting residue elements
to the endomorphism algebra and taking their symmetric
or skew adjoint averages. In symmetric type use
$[q_i]$ and $[q_i a']-[q_i]$; in alternating or paired
type use $H(M_i)$; in unitary type choose a nonzero
residue scalar of the required adjoint sign before
averaging. The sign verification is unchanged because
all adjoints now use the explicitly computed $R$-duality.

For an arbitrary actual form, use the computed direct-sum
isomorphism to the canonical module representatives.
On a self-dual isotypic component its residual Gram matrix is
$\bigl(\pi_i(q_i^{-1}B_{uv})\bigr)_{uv}$, as in (C.3.Q4).
Off-diagonal morphisms between different indecomposable
types lie in the categorical radical, so the matrix is
nonsingular and QSS identifies it with the multiplicity
form. Its determinant or dimension gives its coordinate
in (C.3.P7). Dual-pair multiplicities give the other
integer coordinates. Thus there is no isometry search
left undecided.

\emph{Tensor and exterior operations over $R$.}
Use the computed presentations $M=eR^a$, $N=fR^b$.
Their tensor product is the image of $e\otimes f$
on $R^{ab}$. For $\Lambda^j_RM$ use the image of
$\Lambda^je$ on $\Lambda^jR^a$; its entries are minors.
The coaction on $eR^a$ extends to $R^a$ by putting the
trivial action on $(1-e)R^a$, so its tensor and exterior
coefficient matrices can be computed in free coordinates
and then restricted to these explicit idempotents.
The forms use tensor entries or determinant minors
and similarly restrict to the indicated summands.
The multiplication table of $R$ makes each computation
finite over $k$. If the input degrees are $d,d'$, the output degrees are
at most $d+d'$ and $jd$, respectively. No factorial is divided out.
The images, their $R$-actions and coactions are recovered
by finite $k$-linear algebra, then classified by the
preceding steps.
Finally use bilinearity for virtual products and
$\lambda_t(-x)=\lambda_t(x)^{-1}$ coefficientwise.
Every search in this construction is either finite or
one of the positive algebra-certificate searches with
proved termination. This proves the complete effective
statement with its stated arithmetic access.
\end{proof}

\begin{theorem}[Proved: the effective $C_1$ answer also over imperfect fields]
\label{C:artinian_effective_imperfect_C1}
Let $k$ be any infinite countable $C_1$ field of characteristic
$p>2$, supplied with computable enumeration, exact decidable
arithmetic and equality, and univariate factorization over
every explicitly presented finite tower of field extensions
\[
 k=K_0\subset K_1\subset\cdots\subset K_t,\qquad
             K_{j+1}=K_j[z_j]/(f_j),                     \tag{C.3.P11}
\]
where every $f_j$ is monic irreducible. No perfectness or
separability assumption is made.
For every finite-dimensional commutative $k$-algebra $R$
given by its multiplication table, the complete effective
conclusion of Theorem~\ref{C:artinian_effective_C1} holds:
there are decidable canonical projective module indices,
the explicit form coordinates (C.3.P7), and terminating
algorithms for equality, products and all specified
exterior coefficients. Here the indices are the
natural-number matrix indices, without any Galois-orbit
interpretation.
\end{theorem}
\begin{proof}
We first replace the only perfectness-dependent step in
the effective algebra calculation.
Let $E$ be a finite-dimensional $k$-algebra of dimension $h$.
Enumerate finite products of matrix algebras over the
explicit field towers (C.3.P11), with total $k$-dimension
at most $h$. Each candidate tower is verified to be a
field by the assumed factorization at its successive
stages. Tower coefficients have computable enumerations
in the power bases, so this is a countable effective
search. Enumerate all $k$-linear maps from $E$ to each
candidate product; retain a map exactly when it is
unital, multiplicative, surjective, and its kernel
$J$ has $J^{h+1}=0$. Each test is a finite linear or
ideal-product computation.

A retained kernel is exactly $\operatorname{rad}E$,
since it is nilpotent and the quotient is semisimple.
There is a candidate which is retained. Indeed the
semisimple quotient of $E$ is a product of matrix
algebras over division algebras whose centers are
finite extensions of $k$. All those centers are
$C_1$ by the norm proof of Lemma~\ref{C:C1_field_forms},
including its inseparable case, and have zero Brauer
group. Thus each division algebra is just its center.
Every finite field extension is generated by finitely
many algebraic elements; adjoining them successively
expresses it as a tower (C.3.P11).
This includes finite extensions that are not generated
by a single element. The true semisimple quotient
therefore appears in the search, with nilpotent
radical of exponent at most $h+1$, proving termination.
The polynomial idempotent lift and complementary
recursion from Lemma~\ref{C:C1_algebra_certificates}
remain valid verbatim.

Now follow the complete matrix procedure in the proof
of Theorem~\ref{C:artinian_effective_C1}, using this
tower certificate for every endomorphism algebra.
The finite tuple equations, affine linear projectivity
test, determinant-polynomial isomorphism test and
least-code indexing require only the stated arithmetic
and infinitude of $k$. They have no perfectness
hypothesis. They classify modules directly, so no
geometric descent or separable-center splitting
is being invoked.
The $R$-dual, its antipode precomposition, the form
equations, and the finite-projective tensor and
exterior constructions are equally unchanged.

For a canonical indecomposable, its computed residue
endomorphism field is now given by a finite tower
instead of necessarily one polynomial over $k$.
It still has a computable power basis, exact
arithmetic and equality, and the assumed factorization
access. Square classes are therefore decidable and
have least enumeration representatives.
Any nontrivial adjoint involution makes the field
quadratic over its fixed field, which is again
a finite $C_1$ extension of $k$. Hence the complete
form classification of Lemma~\ref{C:C1_field_forms}
still gives exactly (C.3.P7), including unitary
and alternating types. QSS applies because all
endomorphism algebras remain finite dimensional
with nilpotent radicals. The same residual Gram
extraction consequently computes every coefficient.
This proves the whole algorithm under the tower
access hypothesis. It makes no extension of
the perfect-only Galois-index theorems.
\end{proof}
